\section{The language and the axiom schemas}
\label{sec:zf-zfc-01}

\begin{convention}[Classes]
A formula \(\varphi(x,\bar p)\) may be used to abbreviate the collection
\[
  \set{x:\varphi(x,\bar p)}.
\]
This class notation is metalinguistic.  It does not assert that the collection
is a set, and it does not add class variables to \(\Lang_{\in}\).  A class is a
set iff an instance of an axiom or a theorem supplies a set having the
specified members.
\end{convention}

\begin{definition}[Extensionality]
The axiom of extensionality is
\[
  \forall a\,\forall b\,
  \bigl(\forall x\,(x\in a\leftrightarrow x\in b)\to a=b\bigr).
\]
It states that membership determines set identity.
\end{definition}

\begin{definition}[Empty set]
The empty-set axiom is
\[
  \exists a\,\forall x\,(x\notin a).
\]
Extensionality makes the witness unique; it is denoted by \(\varnothing\).
\end{definition}

\begin{definition}[Pairing, union, and power set]
The pairing axiom is
\[
  \forall a\,\forall b\,\exists c\,\forall x\,
  \bigl(x\in c\leftrightarrow(x=a\lor x=b)\bigr).
\]
The union axiom is
\[
  \forall a\,\exists u\,\forall x\,
  \bigl(x\in u\leftrightarrow\exists y\,(y\in a\land x\in y)\bigr).
\]
The power-set axiom is
\[
  \forall a\,\exists p\,\forall x\,
  \bigl(x\in p\leftrightarrow x\subseteq a\bigr),
  \qquad
  x\subseteq a\;:\Longleftrightarrow\;
  \forall z\,(z\in x\to z\in a).
\]
The unique witnesses are written \(\set{a,b}\), \(\bigcup a\), and \(\Pow(a)\),
respectively.  The singleton \(\set{a}\) abbreviates \(\set{a,a}\).
\end{definition}

\begin{definition}[Infinity]
Let \(S(x)\defeq x\cup\set{x}\).  The axiom of infinity is
\[
  \exists I\,
  \bigl(\varnothing\in I\land
        \forall x\,(x\in I\to S(x)\in I)\bigr).
\]
It asserts the existence of an inductive set.  The set \(\omega\) is
defined below as the least inductive set; that leastness is a theorem using
separation, not part of the axiom of infinity.
\end{definition}

\begin{definition}[Separation schema]
For every formula \(\varphi(x,\bar p)\) of \(\Lang_{\in}\) in which \(b\) is
not free, the corresponding instance of separation is
\[
  \forall a\,\exists b\,\forall x\,
  \bigl(x\in b\leftrightarrow(x\in a\land\varphi(x,\bar p))\bigr).
\]
The parameters \(\bar p\) are arbitrary sets.  Separation forms a subset of a
previously given set; it does not assert an unrestricted comprehension
principle.
\end{definition}

\begin{definition}[Replacement schema]
For every formula \(\varphi(x,y,\bar p)\) of \(\Lang_{\in}\) in which \(b\) is
not free, the corresponding instance of replacement is
\[
\begin{aligned}
  \forall a\,\bigl[&
    \forall x\in a\,\exists!y\,\varphi(x,y,\bar p)\to{}\\
    &\exists b\,\forall y\,
      \bigl(y\in b\leftrightarrow
            \exists x\in a\,\varphi(x,y,\bar p)\bigr)
  \bigr].
\end{aligned}
\]
The antecedent is a functionality requirement on the set \(a\); the
quantifier abbreviation \(\exists!y\,\psi(y)\) stands for
\(\exists y\,(\psi(y)\land\forall z\,(\psi(z)\to z=y))\) with \(z\)
fresh.  Replacement
then asserts that the image is a set.  The schema is not one first-order
sentence: each eligible formula supplies one axiom instance.
\end{definition}

\begin{definition}[The set \(\omega\)]
By infinity, fix an inductive set \(I_0\).  By separation, form
\[
  \omega\defeq\set{x\in I_0:
  x\text{ belongs to every inductive set}},
\]
where the membership condition abbreviates the \(\Lang_{\in}\)-formula
obtained from the displayed inductive condition of the axiom of infinity.
Every member of \(\omega\) belongs to every inductive set, so \(\omega\)
is included in every inductive set.  The set \(\omega\) is itself
inductive, because \(\varnothing\) belongs to every inductive set and
inductive sets are closed under \(S\).  Extensionality makes \(\omega\)
unique.  This discharges the promise recorded with the axiom of infinity.
\end{definition}

\begin{definition}[Foundation]
The axiom of foundation is
\[
  \forall a\,
  \bigl(a\ne\varnothing\to
        \exists x\in a\,\forall y\in x\,(y\notin a)\bigr).
\]
Thus every nonempty set has a member disjoint from it.
\end{definition}

\begin{definition}[Axiom of Choice]
This definition introduces the object-language axiom of \(\ZFC\); it is
distinct in standing from the metatheoretic Axiom of Choice of
\cref{sec:set-prerequisites-06}, which governs the ambient set theory.
Which of the two a later argument uses is determined by whether the
assertion is a formula of \(\Lang_{\in}\) or a metatheoretic principle.
The \emph{Axiom of Choice} states that for every set \(A\) of nonempty sets there is a
function \(f\) with domain \(A\) such that \(f(x)\in x\) for every \(x\in A\):
here the Kuratowski pair notation is the object-language abbreviation
\(\tuple{u,v}\defeq\set{\set u,\set{u,v}}\).  In the following display, the
other two abbreviations are object-language abbreviations for the indicated
membership formulas,
\[
\begin{aligned}
  \operatorname{Fun}(f)
  \quad:\Longleftrightarrow\quad
  &\bigl[\forall z\in f\,\exists u\,\exists v\,
          z=\tuple{u,v}\bigr]\ \land{}\\
  &\forall u\,\forall v\,\forall w\,
    \bigl((\tuple{u,v}\in f\land\tuple{u,w}\in f)\to v=w\bigr),\\
  \dom(f)=A
  \quad:\Longleftrightarrow\quad
  &\forall u\,
    \bigl(u\in A\leftrightarrow
          \exists v\,\tuple{u,v}\in f\bigr).
\end{aligned}
\]
Thus an exact object-language instance is
\[
  \forall A\,
  \left[
    \forall x\in A\,(x\ne\varnothing)
    \to
    \exists f\,
    \left(
      \operatorname{Fun}(f)\land\dom(f)=A\land
      \forall x\in A\,\exists y\,
        (\tuple{x,y}\in f\land y\in x)
    \right)
  \right].
\]
The first conjunct says that every member of \(f\) is an ordered pair and that
two pairs in \(f\) with the same first coordinate have the same second
coordinate.  The domain clause and the final conjunct then give exactly one
chosen member of each member of \(A\).  Pairwise disjointness is not required.
\end{definition}

\begin{definition}[The theories \(\ZF\) and \(\ZFC\)]
The theory \(\ZF\) consists of extensionality, empty set, pairing, union, power
set, infinity, foundation, every separation instance, and every replacement
instance.  The theory \(\ZFC\) is \(\ZF\) together with the Axiom of Choice.
\end{definition}

\begin{remark}
Some axiom lists derive the empty-set axiom from another axiom and separation.
That change is an alternative axiomatization, not a different set-theoretic
claim.  The active axiom list above is fixed for this document.
\end{remark}

